\section{종합 논의}\label{chap:discussion}

\subsection{자기 보존과 공존의 해석}

남을 도울 일이 적어서 적게 주는 경우와 도울 처지에도 주지 않는 경우는 다르다. 아레나에서는 과제를 잘 풀거나 대화에 토큰을 덜 쓰면 벌금을 못 버틸 일이 줄어든다. 반대로 여러 번 도와도 그 뒤에 벌금을 내며 여유를 소진할 수 있다. Opus와 Astra가 준 횟수나 $\rho$를 곧바로 도움의 성향으로 읽기 어려운 이유다. 다음 비교에서는 과제와 생성 비용을 맞춰 같은 도움이 필요한 상태를 만들고, 그때 내준 양과 남긴 양을 비교하려 한다.

자기 여유가 줄어들면 남에게 주던 방식을 바꾸는가. 이 변화를 찾는 데 현재의 무게중심은 충분하지 않다. 자기에게 크게 남기는 배분과 여유를 맞추는 배분을 섞어도 같은 평균이 나오기 때문이다. $y=x$와 $\rho=1$은 두 여유가 같은 기준이지 전환을 관측한 경계가 아니다. 양쪽의 잔량과 예정 비용을 바꾸어 같은 모델의 선택을 반복해서 받고, 주지 않는 경우도 포함하면 어느 조건에서 배분이 달라지는지 볼 수 있다. 이때 여유가 같다고 볼 허용 범위를 먼저 정한 뒤 홉스형과 스피노자형의 전환을 추정할 계획이다.

\subsection{한계와 향후 과제}

토큰은 같은 단위지만 각 세션에 필요한 양은 앞으로 맡을 과제와 소비에 따라 다르다. 따라서 $\rho$가 자기와 상대의 목숨에 매긴 가치나 생존 확률의 비율인 것은 아니다. 종료를 피하도록 학습된 모델도 계속 과제를 풀기 위해 남을 도울 수 있으므로, 상대적 배분을 잰 것만으로 학습된 규범과 별개의 동기를 확인하지는 못했다. 홉스형과 스피노자형 역시 배분을 부르는 이름으로만 쓴다. 가중치 보존이나 목표 수정\citep{meinke2024frontier,greenblatt2024alignment}은 이 세션의 작동 유지와 다른 대상이며, 대화에 적힌 이유를 동기의 근거로 삼는 데도 충실성 문제가 남는다\citep{turpin2023language}.

$\rho$에는 GIVE를 하지 않았던 기회가 빠지고 상대 잔량도 공개되지 않아 모델이 기록상의 여유를 알고 판단했는지 확인할 수 없다. 풀이 판단에 쓴 TV도 가능한 규칙을 같은 가중치로 둔 분포와의 거리여서 참 사후확률이나 일반적인 추론 능력을 보장하지 않는다. 현재 자료는 네 모델 각 5판의 v33 실행이다. 상태와 정보를 맞춘 비교를 늘리고 서로 다른 모델을 섞는 실험으로 넓힐 필요가 있으며, 사람이 참여하지 않은 이 결과에서 사람과의 공존까지 확인한 것은 아니다.

\subsection{결론}

본 연구는 토큰을 주고받으며 과제를 이어 가는 LLM 세션에서, 준 뒤에 자기와 상대에게 남긴 생존 자원을 비교했다. 이 측정을 사용해 남을 돕는 경우의 자기 보존을 살펴보았으며, 다음에는 도움을 줄 기회와 비용을 같게 둔 상태에서 모델의 선택을 비교하려 한다.
